\documentclass[11pt]{article}
\usepackage{preamble}
\renewcommand{\answer}[1]{}

\begin{document}

\ladderheading{AP Statistics \textperiodcentered\ Unit 3 \textperiodcentered\ Topic 3.7 \textperiodcentered\ B: 2026-12-07 \textperiodcentered\ E: 2027-01-13}{Sorting Trays After New Labels}

\focusbanner{TODAY'S PROBLEM}

A cafeteria served 5{,}000 September trays after new sorting-bin labels were installed. The director claims at least $60\%$ are sorted correctly. An SRS of 250 trays finds 136 correct; the club tests at $\alpha=0.05$.\par\smallskip \textbf{Imani:} ``The $p$-value is 0.0354, so reject $H_0$. The labels caused sorting below 60 percent for every month.''\par \textbf{Drew:} ``The 5.6-point gap may be too small to matter, so it is not significant. Accept the director's claim.''\par\smallskip The test conditions have been checked: random sample; $250<500$; null-expected counts 150 and 100.\par
\begin{center}
\textbf{September study}\par\smallskip
\begin{tabular}{|l|c|c|c|}
\hline
\textbf{Source} & \textbf{Trays} & \textbf{Correct} & \textbf{$\alpha$} \\ \hline
\textbf{September} & 5,000 & claim: $\geq 60\%$ & -- \\ \hline
\textbf{SRS} & 250 & 136 & 0.05 \\ \hline
\end{tabular}
\end{center}
\begin{firsttakebox}{FIRST TAKE --- before the video (5 min)}
Can these September trays tell us what caused the result in every month? Give one reason from the study.
\workspace{0.9in}
\end{firsttakebox}

\begin{callout}{\IconBook\ MAKE THE DECISION, THEN BOUND THE CLAIM (keep on the page)}{calloutgreen}
For a one-proportion test, $z=(\widehat p-p_0)/\sqrt{p_0(1-p_0)/n}$; for $H_a:p<p_0$, use the left normal tail.
If the $p$-value is at most the predetermined $\alpha$, reject $H_0$; otherwise fail to reject it.
State evidence for $H_a$ in context. Failing to reject does not establish that $H_0$ is true.
Statistical significance uses the evidence rule; practical importance asks whether an effect is large enough to matter.
A random sample supports inference to its population. A cause requires a comparison with random assignment.
\end{callout}

\newpage
\textbf{Finish after the video:}\par\smallskip

\dokbadge{1}~\textbf{(a)} Define $p$ and state the hypotheses for checking whether correct sorting is below 60\%.\par
\workspace{0.8in}

\dokbadge{2}~\textbf{(b)} Calculate $\widehat p$, the one-proportion $z$ statistic, and the lower-tail $p$-value. Use the verified conditions above.\par
\workspace{1.4in}

\dokbadge{3}~\textbf{(c)} In 4--5 sentences, correct both reports. Compare the $p$-value with $\alpha$, separate statistical significance from practical importance, and limit the conclusion to the people or items studied and what the design can show about cause.\par
\workspace{2.0in}

\begin{sentenceframebox}\raggedright\textbf{Frame:} Because $p\text{-value}=\blank[0.6in]$ is $\blank[0.5in]$ $\alpha$, I would $\blank[0.8in]$ $H_0$ and conclude \blank[2.2in].\end{sentenceframebox}

\begin{sentenceframebox}\raggedright\textbf{Frame:} The result is statistically \blank[0.8in], but the design does not show \blank[2.2in]; the 5.6-point gap's practical meaning depends on \blank[1.5in].\end{sentenceframebox}

\begin{summaryexitbox}{EXIT --- turn this in}
One thing the video changed about my first take: \hrulefill\par
\workspace{0.4in}
\end{summaryexitbox}

\end{document}
